\documentclass{beamer}
\usetheme{Luebeck}
\usecolortheme{beaver}
\usepackage[czech]{babel}
\usepackage[utf8]{inputenc}
\usepackage{color}
\usepackage{graphicx}
\usepackage{amsfonts}
\usepackage{picture}

%\slideCaption{Projekt 5. do předmětu ITY, Vojtěch Večeřa, \today}

\begin{document}

\title{Rasismus}
\subtitle{není jen o barvě pleti\dots} 
%\author{Vojtěch Večeřa}
\author[Vojtěch Večeřa]{Vojtěch Večeřa\\\url{xvecer18@stud.fit.vutbr.cz}}
%\email{xvecer18@stud.fit.vutbr.cz}
\date{\today}
\frame{\titlepage}

\begin{frame}
\transwipe
\transduration{1.5}
\frametitle{Přehled témat}
\begin{itemize}
\item definice rasismu
\item nejrozšířenější typy
\item příklady ze součastnosti
\end{itemize}
\end{frame}


\begin{frame}
\transwipe
\transduration{1.5}
\frametitle{Jak definovat rasismus?}
\begin{columns}
\begin{column}{0.6\textwidth}
\begin{itemize}
\onslide<1->{
\item legálně}
\onslide<2->{
\item sociologicky}
\onslide<3->{
\bigskip
\item jako:
\begin{itemize}
\item xenofobii
\item nadřazenost bílých
\item rozlišování práv různých skupin
\end{itemize}}
\end{itemize}
\end{column}

\begin{column}{0.4\textwidth}
\onslide<1->{
\begin{figure}
\includegraphics[scale=0.5]{prejudice_socialpower.eps}
\end{figure}}
\end{column}
\end{columns}
\end{frame}

\begin{frame}
\transwipe
\transduration{1.5}
\frametitle{Typy rasismu}
\begin{columns}[onlytextwidth]
\begin{column}{0.4\textwidth}
\begin{itemize}
\item rasová diskriminace
\item institucionální
\item ekonomický
\item kulturní
\item barvoslepost
\item vyčleňování
\item nevědomé
\end{itemize}
\end{column}
\begin{column}{0.6\textwidth}
\begin{figure}
\includegraphics[scale=0.56]{colorblidness_m.eps}
\end{figure}
\end{column}
\end{columns}
\end{frame}

\begin{frame}
\transwipe
\transduration{1.5}
\frametitle{Mezimenšinový rasismus}

\quad --\quad chování na základě předsudků mezi menšinami
\newline

\textbf{Převážně}:

\begin{itemize}
\item v Evropě
\item v severní Americe
\end{itemize}
\end{frame}

\begin{frame}
\transwipe
\transduration{1.5}
\frametitle{Evropa}
\onslide<1->{
\textbf{Anglie}

\begin{itemize}
\item Birmingham - černoši proti jiho--asiatům
\item Dewsbury - kurdi proti jiho--asiatům
\end{itemize}}
\onslide<2->{
\textbf{Francie} - 6 milionů muslimů proti 600 tisícům židů

\begin{itemize}
\item převážně muslimé s předky z Arábie
\item nově i Karibičané proti židům
\end{itemize}}
\end{frame}

\begin{frame}
\transwipe
\transduration{1.5}
\frametitle{Severní Amerika}

\textbf{Konflikty mezi}:

\begin{itemize}
\item Afro--američany a Koreo--američany (L.A.)
\item černochy a židy
\item novými imigranty (z latinské Ameriky apod.)
\item agrese vůči bělochům
\item americkými černochy a imigranty z Afriky
\end{itemize}
\end{frame}

\begin{frame}
\transwipe
\transduration{1.5}
\frametitle{Etnická rozmanitost světa}
\begin{figure}[h]
\centering
\includegraphics[scale=0.35]{diverity-map-harvard2.eps}
\end{figure}
\end{frame}

\begin{frame}
\transwipe
\transduration{1.5}
\frametitle{Otázky?}
\begin{columns}
\begin{column}{0.5\textwidth}
\textbf{Věděli jste, že?}

\setlength{\leftskip}{1.3em} 21. březen je mezinárodní den eliminace rasové diskriminace?
\end{column}
\begin{column}{0.5\textwidth}
\begin{figure}[h]
\centering
\includegraphics[scale=0.5]{panda_m.eps}

\end{figure}
\end{column}
\end{columns}
\end{frame}

\begin{frame}
\transwipe
\transduration{1.5}
\frametitle{Zdroje}
\begin{itemize}
\item http://en.wikipedia.org/wiki/Racism
\item http://feminspire.com
\item https://philebersole.wordpress.com
\end{itemize}
\end{frame}
\end{document}